% ---------------------------------------------------------------------------
% Study report · Predicting the inclination from the extended sensor package
%
% Build:  pdflatex inclination_prediction_report.tex   (twice, for the ToC)
% Figures are read from ../outputs/ and are produced by the paired notebook.
% ---------------------------------------------------------------------------
\documentclass[11pt,a4paper]{article}

\newcommand{\runninghead}{Predicting the inclination · station 02 · extended
sensor package}

% ---------------------------------------------------------------------------
% reportstyle.tex — shared preamble for the study reports under studies/.
%
% Kept deliberately inside the package set shipped with TeX Live "basic":
% no tcolorbox, no siunitx, no enumitem, no titlesec. Everything below
% compiles with a stock pdflatex installation.
%
% Usage, from a report directory one level below a study folder:
%     % ---------------------------------------------------------------------------
% reportstyle.tex — shared preamble for the study reports under studies/.
%
% Kept deliberately inside the package set shipped with TeX Live "basic":
% no tcolorbox, no siunitx, no enumitem, no titlesec. Everything below
% compiles with a stock pdflatex installation.
%
% Usage, from a report directory one level below a study folder:
%     % ---------------------------------------------------------------------------
% reportstyle.tex — shared preamble for the study reports under studies/.
%
% Kept deliberately inside the package set shipped with TeX Live "basic":
% no tcolorbox, no siunitx, no enumitem, no titlesec. Everything below
% compiles with a stock pdflatex installation.
%
% Usage, from a report directory one level below a study folder:
%     \input{../../_shared/reportstyle.tex}
%     \graphicspath{{../outputs/}}
% ---------------------------------------------------------------------------

\usepackage[T1]{fontenc}
\usepackage[utf8]{inputenc}
\usepackage{textcomp}
\usepackage{lmodern}
\usepackage{microtype}
\usepackage[margin=2.4cm,headheight=15pt]{geometry}
\usepackage{graphicx}
\usepackage{booktabs}
\usepackage{tabularx}
\usepackage{array}
\usepackage{amsmath}
\usepackage{xcolor}
\usepackage[font=small,labelfont=bf,justification=justified]{caption}
\usepackage{float}
\usepackage{fancyhdr}
\usepackage[hidelinks]{hyperref}

% --- Palette ---------------------------------------------------------------
\definecolor{rulegrey}{HTML}{B9C0C7}
\definecolor{fillgrey}{HTML}{F2F4F6}
\definecolor{failred}{HTML}{A5202B}
\definecolor{passgreen}{HTML}{1B6B3A}
\definecolor{accent}{HTML}{1F4E79}

% --- Semantic shorthands ---------------------------------------------------
\newcommand{\degC}{\textdegree C}
\newcommand{\mdegC}{mdeg/\textdegree C}
\newcommand{\fail}{\textcolor{failred}{\textbf{fail}}}
\newcommand{\pass}{\textcolor{passgreen}{\textbf{pass}}}
\newcommand{\worse}{\textcolor{failred}{worse}}
\newcommand{\better}{\textcolor{passgreen}{better}}

% --- Highlighted panel -----------------------------------------------------
% #1 = heading, #2 = body
\newcommand{\panel}[2]{%
  \begin{center}%
  \setlength{\fboxsep}{9pt}%
  \fcolorbox{rulegrey}{fillgrey}{%
    \begin{minipage}{0.93\textwidth}%
      {\bfseries\color{accent}#1}\par\medskip
      #2
    \end{minipage}}%
  \end{center}}

% --- Numeric column: right aligned, fixed width ----------------------------
\newcolumntype{R}[1]{>{\raggedleft\arraybackslash}p{#1}}
\newcolumntype{L}[1]{>{\raggedright\arraybackslash}p{#1}}
\newcolumntype{Y}{>{\raggedright\arraybackslash}X}

% --- Table look ------------------------------------------------------------
\renewcommand{\arraystretch}{1.15}
\setlength{\tabcolsep}{6pt}

% --- Headers ---------------------------------------------------------------
\pagestyle{fancy}
\fancyhf{}
\fancyhead[L]{\small\itshape\runninghead}
\fancyhead[R]{\small\thepage}
\renewcommand{\headrulewidth}{0.4pt}

% --- Section spacing without titlesec --------------------------------------
\makeatletter
\renewcommand\section{\@startsection{section}{1}{\z@}%
  {-3.2ex \@plus -1ex \@minus -.2ex}{2.0ex \@plus.2ex}%
  {\normalfont\Large\bfseries\color{accent}}}
\renewcommand\subsection{\@startsection{subsection}{2}{\z@}%
  {-2.6ex \@plus -1ex \@minus -.2ex}{1.4ex \@plus.2ex}%
  {\normalfont\large\bfseries}}
\makeatother

% --- Figure defaults -------------------------------------------------------
\setkeys{Gin}{width=\linewidth}
\captionsetup[figure]{skip=6pt}

%     \graphicspath{{../outputs/}}
% ---------------------------------------------------------------------------

\usepackage[T1]{fontenc}
\usepackage[utf8]{inputenc}
\usepackage{textcomp}
\usepackage{lmodern}
\usepackage{microtype}
\usepackage[margin=2.4cm,headheight=15pt]{geometry}
\usepackage{graphicx}
\usepackage{booktabs}
\usepackage{tabularx}
\usepackage{array}
\usepackage{amsmath}
\usepackage{xcolor}
\usepackage[font=small,labelfont=bf,justification=justified]{caption}
\usepackage{float}
\usepackage{fancyhdr}
\usepackage[hidelinks]{hyperref}

% --- Palette ---------------------------------------------------------------
\definecolor{rulegrey}{HTML}{B9C0C7}
\definecolor{fillgrey}{HTML}{F2F4F6}
\definecolor{failred}{HTML}{A5202B}
\definecolor{passgreen}{HTML}{1B6B3A}
\definecolor{accent}{HTML}{1F4E79}

% --- Semantic shorthands ---------------------------------------------------
\newcommand{\degC}{\textdegree C}
\newcommand{\mdegC}{mdeg/\textdegree C}
\newcommand{\fail}{\textcolor{failred}{\textbf{fail}}}
\newcommand{\pass}{\textcolor{passgreen}{\textbf{pass}}}
\newcommand{\worse}{\textcolor{failred}{worse}}
\newcommand{\better}{\textcolor{passgreen}{better}}

% --- Highlighted panel -----------------------------------------------------
% #1 = heading, #2 = body
\newcommand{\panel}[2]{%
  \begin{center}%
  \setlength{\fboxsep}{9pt}%
  \fcolorbox{rulegrey}{fillgrey}{%
    \begin{minipage}{0.93\textwidth}%
      {\bfseries\color{accent}#1}\par\medskip
      #2
    \end{minipage}}%
  \end{center}}

% --- Numeric column: right aligned, fixed width ----------------------------
\newcolumntype{R}[1]{>{\raggedleft\arraybackslash}p{#1}}
\newcolumntype{L}[1]{>{\raggedright\arraybackslash}p{#1}}
\newcolumntype{Y}{>{\raggedright\arraybackslash}X}

% --- Table look ------------------------------------------------------------
\renewcommand{\arraystretch}{1.15}
\setlength{\tabcolsep}{6pt}

% --- Headers ---------------------------------------------------------------
\pagestyle{fancy}
\fancyhf{}
\fancyhead[L]{\small\itshape\runninghead}
\fancyhead[R]{\small\thepage}
\renewcommand{\headrulewidth}{0.4pt}

% --- Section spacing without titlesec --------------------------------------
\makeatletter
\renewcommand\section{\@startsection{section}{1}{\z@}%
  {-3.2ex \@plus -1ex \@minus -.2ex}{2.0ex \@plus.2ex}%
  {\normalfont\Large\bfseries\color{accent}}}
\renewcommand\subsection{\@startsection{subsection}{2}{\z@}%
  {-2.6ex \@plus -1ex \@minus -.2ex}{1.4ex \@plus.2ex}%
  {\normalfont\large\bfseries}}
\makeatother

% --- Figure defaults -------------------------------------------------------
\setkeys{Gin}{width=\linewidth}
\captionsetup[figure]{skip=6pt}

%     \graphicspath{{../outputs/}}
% ---------------------------------------------------------------------------

\usepackage[T1]{fontenc}
\usepackage[utf8]{inputenc}
\usepackage{textcomp}
\usepackage{lmodern}
\usepackage{microtype}
\usepackage[margin=2.4cm,headheight=15pt]{geometry}
\usepackage{graphicx}
\usepackage{booktabs}
\usepackage{tabularx}
\usepackage{array}
\usepackage{amsmath}
\usepackage{xcolor}
\usepackage[font=small,labelfont=bf,justification=justified]{caption}
\usepackage{float}
\usepackage{fancyhdr}
\usepackage[hidelinks]{hyperref}

% --- Palette ---------------------------------------------------------------
\definecolor{rulegrey}{HTML}{B9C0C7}
\definecolor{fillgrey}{HTML}{F2F4F6}
\definecolor{failred}{HTML}{A5202B}
\definecolor{passgreen}{HTML}{1B6B3A}
\definecolor{accent}{HTML}{1F4E79}

% --- Semantic shorthands ---------------------------------------------------
\newcommand{\degC}{\textdegree C}
\newcommand{\mdegC}{mdeg/\textdegree C}
\newcommand{\fail}{\textcolor{failred}{\textbf{fail}}}
\newcommand{\pass}{\textcolor{passgreen}{\textbf{pass}}}
\newcommand{\worse}{\textcolor{failred}{worse}}
\newcommand{\better}{\textcolor{passgreen}{better}}

% --- Highlighted panel -----------------------------------------------------
% #1 = heading, #2 = body
\newcommand{\panel}[2]{%
  \begin{center}%
  \setlength{\fboxsep}{9pt}%
  \fcolorbox{rulegrey}{fillgrey}{%
    \begin{minipage}{0.93\textwidth}%
      {\bfseries\color{accent}#1}\par\medskip
      #2
    \end{minipage}}%
  \end{center}}

% --- Numeric column: right aligned, fixed width ----------------------------
\newcolumntype{R}[1]{>{\raggedleft\arraybackslash}p{#1}}
\newcolumntype{L}[1]{>{\raggedright\arraybackslash}p{#1}}
\newcolumntype{Y}{>{\raggedright\arraybackslash}X}

% --- Table look ------------------------------------------------------------
\renewcommand{\arraystretch}{1.15}
\setlength{\tabcolsep}{6pt}

% --- Headers ---------------------------------------------------------------
\pagestyle{fancy}
\fancyhf{}
\fancyhead[L]{\small\itshape\runninghead}
\fancyhead[R]{\small\thepage}
\renewcommand{\headrulewidth}{0.4pt}

% --- Section spacing without titlesec --------------------------------------
\makeatletter
\renewcommand\section{\@startsection{section}{1}{\z@}%
  {-3.2ex \@plus -1ex \@minus -.2ex}{2.0ex \@plus.2ex}%
  {\normalfont\Large\bfseries\color{accent}}}
\renewcommand\subsection{\@startsection{subsection}{2}{\z@}%
  {-2.6ex \@plus -1ex \@minus -.2ex}{1.4ex \@plus.2ex}%
  {\normalfont\large\bfseries}}
\makeatother

% --- Figure defaults -------------------------------------------------------
\setkeys{Gin}{width=\linewidth}
\captionsetup[figure]{skip=6pt}

\graphicspath{{../outputs/}}

\title{\vspace{-1.2cm}\textbf{Predicting the inclination of the Gubbio wall
from the extended sensor package}\\[2mm]
\large Station 02, 21 February 2025 to 10 August 2026, hourly grid}
\author{Study report · \texttt{studies/inclination\_prediction/}}
\date{12 August 2026}

\begin{document}
\maketitle
\thispagestyle{fancy}

\begin{abstract}
\noindent
The instrument package installed at station 02 on 21 February 2025 records air
temperature, relative humidity, wall temperature, solar radiation, battery
voltage and inclination on one logger. This report asks what that package buys
as a \emph{prediction} system. Three questions are answered: which single channel
best predicts the inclination, whether a combination does better and how the best
combination should be diagnosed, and how far ahead NeuralProphet can forecast
with a defensible statement of uncertainty. Every driver is given the chance to
act through a transport delay and a thermal inertia before it is ranked, with the
delay bounded at twelve hours on physical grounds, and the ranking is
out-of-sample throughout.

The headline results are that air temperature dominates every ranking; that on
the long summer block the pair \texttt{tair + sr} reduces the error by 27\,\%
over air temperature alone, while wall temperature appears in no winning subset
in either tier once implausible delays are excluded; that the useful forecast
band is roughly three to twelve hours ahead, outside which the model does no
better than repeating the previous day; and that the nominal 90\,\% prediction
interval of the deployable configuration actually covers 54--67\,\% of
observations, so it must not be quoted as ninety.

One structural finding governs the reading of all of it. The site predicts a
\emph{negative} association between the inclination and any heating driver: the
valley face is the more sun-exposed, expands further, and tips the wall towards
the mountain. The compensated target does show that sign, at $r = -0.956$ with
air temperature. The \emph{raw} channel does not --- it moves \emph{positively},
at $r = +0.845$ and $+1.63$\,mdeg\,/\,\degC. The expected sign in the target is
produced by the compensation subtracting 5\,\mdegC\ from a channel whose measured
slope is 1.63, not by the structure.

A separate diagnostic bears on the same question. Given a signed operator, the
inclination is found to \emph{lead} the wall-temperature probe by two hours, on
the raw channel as well as the compensated one --- the signature of a probe set
deeper than the material whose expansion drives the movement. The probe's depth
is not recorded anywhere in the archive; the phase measured here is the check to
run against it once it is.
\end{abstract}

\tableofcontents

% ===========================================================================
\section{What this study asks, and what it assumes}

Three questions, posed of the extended sensor package:

\begin{enumerate}
\item Over the period in which the extended set of sensors exists, is it
possible to use one of the sensor variables to predict the inclination, and
which is the best predictor?
\item Can a set of more than one predictor perform better than one alone, and
how should the best combination be diagnosed?
\item Using the best scenario, how far into the future can the inclination be
predicted, and with what measure of uncertainty, using NeuralProphet?
\end{enumerate}

\subsection{Premises, imposed rather than tested}

Three premises are fixed by the study design. They differ deliberately from
those of the companion study \texttt{studies/thermal\_compensation/}, which
examined exactly these assumptions and rejected two of them. This study accepts
them and asks a different question on top.

\begin{center}
\begin{tabularx}{\textwidth}{L{4.6cm}Y}
\toprule
\textbf{Premise} & \textbf{Consequence for this report} \\
\midrule
The logged \texttt{I} channel is an inclination in millidegrees, not a
conditioner voltage &
The 2500 offset that the companion study subtracts is added back at load time.
Nothing downstream is affected: the compensation normalises the series to start
at zero, which removes every additive constant. Only the printed levels change.
\\
The documented temperature compensation is correct &
The compensated series \texttt{inc\_comp} is the target of every model below.
The coefficient is never swept and the acceptance tests of the companion study
are not repeated. \\
The analysis grid is one hour &
Every external proxy the production pipeline aligns against is published hourly,
so the record is aggregated immediately after parsing and every correlation,
operator scan, fit and forecast in this report is computed on that grid. \\
\bottomrule
\end{tabularx}
\end{center}

\subsection{The site's expected sign}
\label{sec:sign}

The wall stands in an embankment situation with two distinct faces. The
\textbf{valley side} receives markedly more direct insolation than the
\textbf{mountain side}, so under daytime heating it expands further, and the
differential expansion tips the wall \textbf{towards the mountain}. Under the
instrument's convention that is a \textbf{negative} change in inclination.

The site physics therefore predicts a \emph{negative} association between the
inclination and every heating driver --- air temperature, wall temperature,
solar radiation. A negative correlation is the expected result here and must not
be reported as an anomaly or corrected away.

This is a prediction about the \emph{structure}, so it has to be tested on the
signal before the compensation touches it. The distinction is not academic: the
compensation subtracts 5\,\mdegC, and if the instrument's true thermal slope is
smaller than that, the residual carries the opposite sign to the original by
arithmetic alone. A negative sign already present in the raw channel is
structural; a negative sign that appears only after compensation is not.

\subsection{What the raw channel actually shows}
\label{sec:signtest}

\begin{center}
\small
\begin{tabular}{llR{1.4cm}R{1.7cm}R{1.4cm}R{1.7cm}}
\toprule
\textbf{Series} & \textbf{Driver} & \textbf{$r$ lev.} &
\textbf{slope lev.} & \textbf{$r$ band} & \textbf{slope band} \\
\midrule
Raw \texttt{inc}       & \texttt{tair}  & $+0.845$ & $+1.628$ & $+0.966$ & $+2.291$ \\
                       & \texttt{twall} & $+0.791$ & $+1.248$ & $+0.767$ & $+1.319$ \\
                       & \texttt{sr}    & $+0.588$ & $+0.021$ & $+0.758$ & $+0.017$ \\
\addlinespace
\texttt{inc\_comp}     & \texttt{tair}  & $-0.956$ & $-3.372$ & $-0.975$ & $-2.709$ \\
                       & \texttt{twall} & $-0.878$ & $-2.535$ & $-0.701$ & $-1.411$ \\
                       & \texttt{sr}    & $-0.478$ & $-0.032$ & $-0.803$ & $-0.022$ \\
\bottomrule
\end{tabular}
\end{center}

\noindent{\footnotesize Tier~1 window. Slopes are in mdeg per \degC\ for the
temperatures and mdeg per W/m$^2$ for radiation; ``band'' denotes the diurnal
band, with a one-week centred rolling mean removed.\par}

\panel{The expected sign is present in the target, and absent from the
instrument}{%
The raw inclination moves \emph{positively} with every heating driver: $+0.845$
with air temperature on levels, $+0.966$ on the diurnal band, at a slope of
$+1.63$\,\mdegC. This is the opposite of what the embankment geometry predicts.

\medskip
The compensated series shows the expected negative sign --- but it does so by
arithmetic. The compensation subtracts 5\,\mdegC\ from a channel whose measured
slope is $+1.63$, leaving a residual slope of $1.63 - 5 = -3.37$\,\mdegC, which
is exactly the value the table reports. The sign flip is entirely a product of
the over-correction.

\medskip
Three readings are consistent with this, and these data cannot separate them.
\textbf{(i)} The instrument's sign convention is inverted relative to the one
assumed, so a positive raw change \emph{is} movement towards the mountain.
\textbf{(ii)} The instrument's own thermal drift is positive and larger than the
structural response, masking it. \textbf{(iii)} The structural response at this
station does not follow the differential-expansion expectation. Settling this
requires the calibration sheet, the installation record, or an independent
measurement of the mounting orientation --- none of which is in the archive. The
contradiction is recorded in \texttt{docs/raw-data-format.md} Section~7.5 rather
than resolved here.}

\subsection{What this costs the rest of the report}

Every statement of the form ``air temperature is the best predictor'' is true of
this target and means, in substantial part, that air temperature predicts the
residual of its own correction. Three consequences are carried throughout.

\textbf{The relative standings survive.} Every configuration is scored against
the same target, so the ordering of channels, the operator estimates, the
combination diagnostics and the calibration results are all internally valid.

\textbf{The absolute errors do not transfer.} They describe a system predicting a
partly artificial signal, and should not be quoted for a raw-inclination
deployment without re-running the study on that target.

\textbf{The perfect-foresight regime of Section~\ref{sec:horizon} is inflated by
the same mechanism.} Handing the model future air temperature is close to handing
it the future target, because the target contains $-5\,T_{\mathrm{air}}$ by
construction. That regime is read accordingly there.

% ===========================================================================
\section{Data, cleaning and the windows the study can use}

\subsection{Loading}

The loader implements the contract in \texttt{docs/raw-data-format.md}: decimal
separators normalised per field rather than per file, record timestamps
preferred over filename dates, duplicate timestamps merged field by field rather
than dropped, and every sentinel mapped to \texttt{NaN} --- zeros in the
battery, temperature, humidity and inclination channels, $-55$ in the wall
probe, and the unsigned wrap-around above 1400\,W/m$^2$ in the pyranometer.
Night-time zeros in solar radiation are kept, because they are measurements.

Hourly aggregation takes the mean of each interval and accepts the interval only
when at least two of its three raw samples survived cleaning, so an hour rebuilt
from a single reading is not presented as equivalent to a complete one. A
rolling-median context filter then removes the instrument-settling transient that
follows the 103-day outage.

\begin{figure}[H]
\includegraphics{S3_F01_st02_overview}
\caption{The compensated inclination and its four candidate drivers over the
whole extended-package era, on the hourly grid. The two long interruptions ---
the wall-probe failure from mid-July 2025 and the acquisition outage from March
to June 2026 --- are visible as the gaps in the \texttt{twall} panel and in every
panel respectively. The inclination's annual excursion is roughly 150\,mdeg
against a diurnal amplitude of a few tens, and its shape is close to an inverted
air-temperature curve, which Section~\ref{sec:signtest} attributes to the
correction rather than to the structure.}
\end{figure}

\begin{figure}[H]
\includegraphics{S3_F02_st02_target}
\caption{Raw and compensated inclination, with the correction term that separates
them. The correction is not a small adjustment: over the era it spans a range
comparable to the signal's own, and the compensated series is close to the
negative of the correction. This is the figure behind the residual slope of
$-3.37$\,\mdegC.}
\end{figure}

\subsection{Which periods carry which channels}

Wall temperature is the binding constraint of this era. The probe reads its
open-circuit sentinel on 219 of the 416 days present, while every other channel
is available whenever the logger is. That asymmetry forces the study into two
tiers, and the choice between them is a substantial part of Question~2.

\begin{center}
\small
\begin{tabular}{L{5.6cm}R{2.2cm}R{1.6cm}R{1.8cm}}
\toprule
\textbf{Longest contiguous runs} & \textbf{Start} & \textbf{Days} &
\textbf{Coverage} \\
\midrule
\multicolumn{4}{l}{\itshape All channels, wall temperature included} \\
\quad spring block            & 2025-02-21 & 73.1  & 99.5\,\% \\
\quad early-summer block      & 2025-05-05 & 69.6  & 99.9\,\% \\
\quad detached summer block   & 2026-06-18 & 53.3  & 95.7\,\% \\
\addlinespace
\multicolumn{4}{l}{\itshape Wall temperature dropped} \\
\quad summer block            & 2025-05-05 & 143.2 & 99.9\,\% \\
\quad winter block            & 2025-11-25 & 101.0 & 99.7\,\% \\
\quad spring block            & 2025-02-21 & 73.1  & 99.5\,\% \\
\quad detached summer block   & 2026-06-18 & 53.3  & 95.7\,\% \\
\bottomrule
\end{tabular}
\end{center}

\textbf{Tier 1} is the calendar window 2025-02-21 to 2025-07-13, 3432 hourly
slots, in which all four drivers are present. Internally it is not one run but
two, split on 5 May by a short wall-probe interruption --- a fact that becomes
important in Section~\ref{sec:horizon}, because it means any autoregressive model
fitted on Tier~1 is fitted across a discontinuity. The detached summer block of
2026 is held back entirely as an out-of-period test.

\textbf{Tier 2} drops wall temperature and takes the longest contiguous run that
remains: 2025-05-05 to 2025-09-25, 143.2 days, 3438 slots at 99.9\,\% coverage.
It is not longer than Tier~1 in days --- the era contains no longer clean run ---
but it is what a deployed system actually has, it extends two and a half months
further into the year, and it is unbroken.

The two tiers overlap between 5 May and 13 July. They are therefore not
independent samples, and agreement between them is weaker evidence than it would
otherwise be. Where they disagree, the disagreement is informative.

\begin{figure}[H]
\includegraphics{S3_F03_st02_diurnal}
\caption{Mean diurnal cycle over the Tier~1 window. The inclination carries a
daily cycle of roughly 20\,mdeg, in anti-phase with air temperature --- peaking
near dawn, troughing in mid-afternoon. That is the sign the embankment geometry
predicts. Section~\ref{sec:signtest} shows that in this target it is produced by
the correction: the same cycle drawn on the raw channel is in phase with air
temperature, not against it.}
\end{figure}

\begin{figure}[H]
\includegraphics{S3_F04_st02_correlations}
\caption{Correlation structure over the Tier~1 window, on levels (left) and on
first differences (right). Levels correlations between two trending series are
inflated by the shared trend; the differenced matrix is the honest one for a
signal with drift. Both are shown so the gap is visible. Note the block of high
mutual correlation among the drivers themselves, quantified in
Section~\ref{sec:collinearity}.}
\end{figure}

% ===========================================================================
\section{Question 1 --- the operator scan and the best single predictor}

\subsection{Why an operator is fitted before anything is ranked}

A structure with thermal mass cannot respond to a change in forcing
instantaneously, so before any driver is ranked it is given the chance to act
through the operator that suits it. Two operators are scanned \emph{jointly},
because they are physically distinct and scanning either alone can attribute to
one what belongs to the other:

\begin{itemize}
\item a \textbf{transport delay} $d$ shifts the driver in time without changing
its shape --- what a thermal front travelling from the exposed face to the depth
that moves the instrument would do;
\item a \textbf{thermal inertia} $\tau$ low-passes the driver through a
single-pole filter with $\alpha = \Delta t/(\tau + \Delta t)$ --- what a body
integrating the forcing would do. The instantaneous case is $\tau = 0$ and
remains in the grid, so it competes rather than being assumed away.
\end{itemize}

\subsection{The delay is bounded at twelve hours}
\label{sec:cap}

The grid is 13 delays (0 to 12 hours) by 16 time constants (0 to 168 hours) for
each driver, scored by explained variance. The bound on the delay is physical and
it is load-bearing.

\textbf{The mechanism sets the scale.} Direct insolation heats one face of a
masonry embankment. A thermal front reaching the depth that governs the
inclination is not expected to take longer than half a day, and no process at
this site would impose a slower transport.

\textbf{Beyond half a cycle the scan stops being interpretable.} On a
near-periodic diurnal signal, a delay of $d$ hours is indistinguishable from a
\emph{lead} of $24 - d$. A grid extending past twelve hours can therefore return
an optimum that is really an aliased lead --- which no causal filter can
represent and no forecasting system can use. Twelve hours is the half-cycle, and
stopping there keeps every reported optimum interpretable.

\panel{The bound is correct for the forcings, and the wrong shape for the wall
probe}{%
This bound is derived for \textbf{external forcings}. Air temperature and solar
radiation are inputs to the system: they must precede the response they cause, so
constraining their delay to be non-negative is physically right, and
Section~\ref{sec:probe} confirms that both of them do in fact peak at a
non-negative delay.

\medskip
\textbf{Wall temperature is not a forcing.} It is an \emph{internal state
variable} --- a temperature measured at some unrecorded depth inside the masonry
--- and nothing requires it to precede the deformation. An earlier version of
this report scanned delays to 48 hours, found an optimum for \texttt{twall} at
\textbf{20 hours}, and dismissed it as lying ``outside the physically admissible
range''. That dismissal was wrong in kind. Twenty hours is the 24-hour alias of a
four-hour \emph{lead}, and a lead of an internal probe over the response is
admissible --- it is what a deep installation produces.

\medskip
The right correction is not a tighter bound but a \emph{signed} one, and
Section~\ref{sec:probe} applies it. What the causal constraint does remain right
about is \textbf{use}: a lead cannot be applied by a forecasting system, because
it would require the probe's future values at the moment the forecast is issued.
So the constraint stays in force for every model in this report, and the lead is
measured separately as a diagnostic. With the constraint applied, wall
temperature's predictive optimum is \textbf{zero delay, zero time constant} at
$R^2 = 0.491$, it ranks second at $R^2 = 0.459$ out of sample, and it appears in
no winning subset in Section~\ref{sec:combinations}.}

\subsection{Both bands are needed}

The scan is run twice: on the series as they are, and on the diurnal band alone,
with a one-week centred rolling mean removed.

\begin{center}
\small
\begin{tabular}{lR{1.5cm}R{1.3cm}R{1.5cm}R{1.6cm}R{2.0cm}}
\toprule
\textbf{Driver} & \textbf{delay} & \textbf{$\tau$} & \textbf{$r$} &
\textbf{$R^2$} & \textbf{$R^2$ at $d=\tau=0$} \\
\midrule
\multicolumn{6}{l}{\itshape Full band, Tier 1} \\
\quad \texttt{tair}  &  0 h &   0 h & $-0.956$ & 0.914 & 0.914 \\
\quad \texttt{twall} & 12 h & 120 h & $-0.888$ & 0.788 & 0.771 \\
\quad \texttt{sr}    &  0 h & 168 h & $-0.813$ & 0.662 & 0.229 \\
\quad \texttt{rh}    &  0 h & 168 h & $+0.853$ & 0.728 & 0.472 \\
\addlinespace
\multicolumn{6}{l}{\itshape Diurnal band, Tier 1 --- carried forward} \\
\quad \texttt{tair}  &  0 h &   0 h & $-0.975$ & 0.951 & 0.951 \\
\quad \texttt{sr}    &  0 h &   1 h & $-0.837$ & 0.700 & 0.644 \\
\quad \texttt{twall} &  0 h &   0 h & $-0.701$ & 0.491 & 0.491 \\
\quad \texttt{rh}    &  0 h &   0 h & $+0.695$ & 0.483 & 0.483 \\
\bottomrule
\end{tabular}
\end{center}

The two bands disagree sharply, and the disagreement is the point. On the full
band the optima for \texttt{sr} and \texttt{rh} both sit at $\tau = 168$\,hours
--- the last value in the grid, that is, at the boundary --- lifting explained
variance for \texttt{sr} from 0.229 to 0.662. A time constant of one week is not
a thermal property of a masonry wall; it is the width of filter needed to turn a
daily radiation cycle into a seasonal envelope. The full-band optimum is fitting
the annual component, and it is pinned at the grid edge because a longer filter
would fit it better still. Wall temperature does the same thing in two dimensions
at once, running to the edge in both delay and $\tau$.

Once the weekly rolling mean is removed, every optimum moves to zero delay and a
short time constant. Reporting the full-band operators as physical constants
would have produced three confident and entirely spurious findings. The
diurnal-band operators are the ones carried forward.

\begin{figure}[H]
\includegraphics{S3_F05_st02_operators_fullband}
\caption{Explained variance over the delay-by-inertia grid, full band. The bright
band along the bottom edge of the \texttt{sr} and \texttt{rh} panels is the
boundary optimum discussed above: the surface is still rising as $\tau$ reaches
the end of the grid, which is the visual signature of a filter fitting a
component slower than anything the grid contains.}
\end{figure}

\begin{figure}[H]
\includegraphics{S3_F06_st02_operators_diurnal}
\caption{The same grid on the diurnal band, with the delay bounded at twelve
hours. Every optimum now sits in the top-left corner --- zero delay, zero or
one-hour time constant. The gradual brightening towards the right-hand edge of
the \texttt{twall} and \texttt{sr} panels is the leading edge of the 24-hour
alias returning; Section~\ref{sec:cap} is the reason the grid stops before it
becomes a maximum.}
\end{figure}

\subsection{No driver needs a transport delay}

The result of bounding the scan is unusually clean: on the diurnal band,
\textbf{every driver's optimum is at zero delay}, and only solar radiation takes
a non-zero time constant, at one hour.

That is a substantive finding, not a null one. Within the physically admissible
window there is no evidence of a lagged structural response at all --- the
compensated inclination tracks its drivers within the hour. It is also consistent
with Section~\ref{sec:signtest}: a target that contains $-5\,T_{\mathrm{air}}$ by
construction must follow air temperature instantaneously, because that term
\emph{is} air temperature. The finding therefore describes the target's
construction at least as much as the structure's thermodynamics, and a repeat on
the raw channel would be needed to separate the two.

\subsection{The wall probe's phase, and what it says about its depth}
\label{sec:probe}

Nothing in this subsection feeds any model. The operators used everywhere else
are the causal ones fixed above; the guard that enforces this is described at the
end of the subsection, and its effect is verified.

\subsubsection*{Why the wall probe deserves a signed scan}

The deformation is governed by the temperature field near the sun-exposed valley
face. The probe sits at some depth inside the masonry, and that depth decides
which way the phase between them runs:

\begin{center}
\small
\begin{tabularx}{\textwidth}{L{2.6cm}L{6.4cm}Y}
\toprule
\textbf{Probe position} & \textbf{What it reads} & \textbf{Expected phase} \\
\midrule
\textbf{Shallow} & a fast, weakly damped signal close to the surface forcing &
the inclination \textbf{lags} the probe \\
\textbf{Deep} & a strongly damped and delayed signal &
the inclination \textbf{leads} the probe \\
\bottomrule
\end{tabularx}
\end{center}

The depth is not recorded anywhere in the archive. So what follows measures the
\emph{phase}, in hours, and leaves the depth open. When the depth is established
it becomes a consistency check against these numbers rather than a reason to
repeat the analysis.

\subsubsection*{The measurement}

Phase is measured on the diurnal band over signed delays from $-12$ to $+12$
hours. A negative phase means the second channel \emph{lags} the first.

\begin{center}
\small
\begin{tabular}{llR{1.8cm}R{1.8cm}R{1.8cm}R{1.6cm}}
\toprule
\textbf{First} & \textbf{Second} & \textbf{phase} & \textbf{$r$ there} &
\textbf{$r$ at zero} & \textbf{$R^2$ gain} \\
\midrule
\multicolumn{6}{l}{\itshape The wall probe against the response} \\
\quad \texttt{inc\_comp} & \texttt{twall} & $-2$ h & $-0.875$ & $-0.701$ & 0.275 \\
\quad \texttt{inc} (raw) & \texttt{twall} & $-2$ h & $+0.866$ & $+0.767$ & 0.162 \\
\addlinespace
\multicolumn{6}{l}{\itshape The wall probe against what heats it} \\
\quad \texttt{twall} & \texttt{tair} & $+2$ h & $+0.897$ & $+0.753$ & 0.238 \\
\quad \texttt{twall} & \texttt{sr}   & $+3$ h & $+0.891$ & $+0.532$ & 0.510 \\
\addlinespace
\multicolumn{6}{l}{\itshape Controls --- external forcings, must not lag} \\
\quad \texttt{inc\_comp} & \texttt{tair} & $0$ h & $-0.975$ & $-0.975$ & 0.000 \\
\quad \texttt{inc\_comp} & \texttt{sr}   & $+1$ h & $-0.805$ & $-0.803$ & 0.004 \\
\bottomrule
\end{tabular}
\end{center}

\begin{figure}[H]
\includegraphics{S3_F15_st02_operators_signed}
\caption{The diurnal-band operator grid re-run over signed delays. Air
temperature, solar radiation and relative humidity all peak at zero delay, as
external forcings must. Wall temperature peaks at $-4$ hours with a two-hour time
constant, and its entire high-$R^2$ region lies in negative delay. This figure is
diagnostic and feeds no model.}
\end{figure}

\begin{center}
\small
\begin{tabular}{lR{1.6cm}R{1.3cm}R{1.6cm}R{1.8cm}R{2.4cm}}
\toprule
\textbf{Driver} & \textbf{delay} & \textbf{$\tau$} & \textbf{$R^2$ signed} &
\textbf{$R^2$ causal} & \textbf{lost to causality} \\
\midrule
\texttt{tair}  &  0 h & 0 h & 0.951 & 0.951 & 0.000 \\
\texttt{twall} & $-4$ h & 2 h & \textbf{0.797} & 0.491 & \textbf{0.306} \\
\texttt{sr}    &  0 h & 1 h & 0.700 & 0.700 & 0.000 \\
\texttt{rh}    &  0 h & 0 h & 0.483 & 0.483 & 0.000 \\
\bottomrule
\end{tabular}
\end{center}

\subsubsection*{What it shows}

\textbf{The controls pass cleanly.} Air temperature, solar radiation and relative
humidity all peak at a non-negative delay and lose \emph{exactly nothing} to the
causal constraint. Had an external forcing appeared to lag the response, every
number here would have been a method artefact. None does.

\textbf{Wall temperature is the sole exception, and the effect is large.} Its
optimum sits at $-4$ hours with a two-hour time constant, at $R^2 = 0.797$
against 0.491 under the causal constraint --- \textbf{0.306 of explained variance
suppressed}, a 62\,\% relative loss. Nothing else in the package loses anything.

\textbf{The lead survives on the raw channel.} This is the load-bearing check.
The compensated and the raw inclination both lead \texttt{twall} by two hours, at
$-0.875$ and $+0.866$ respectively --- opposite signs, as
Section~\ref{sec:signtest} requires, but the \emph{same phase}. The lead is
therefore a property of the instrument pair, not an artefact of the correction.

\textbf{The chain is internally consistent.} The probe lags solar radiation by
three hours and air temperature by two, while the inclination lags air
temperature by none. Everything is consistent with one picture: the deformation
tracks the near-surface thermal state essentially in step with the air, and the
probe registers the same event two to four hours later because it is set further
in.

\panel{Read against the depth: the probe is deeper than the layer that moves the
wall}{%
The inclination \emph{leads} the wall probe. On the table above, that is the
signature of a \textbf{deep} installation: the probe is set behind the material
whose expansion actually drives the movement, so it reports the thermal event
after the structure has already responded to it.

\medskip
This report deliberately stops short of converting the phase into a depth in
centimetres, which would require assuming a thermal diffusivity for the masonry.
The phase itself is the durable measurement. When the installation depth is
established, the check is direct: a deeper probe should show a larger lag against
air temperature than the two hours measured here, and a correspondingly larger
lead of the inclination over it.

\medskip
There is a corollary that bears on the unresolved sign contradiction of
Section~\ref{sec:signtest}. A deformation genuinely driven by bulk thermal
expansion of masonry ought to lag air temperature by \emph{something}. This one
does not --- it is in step with the air to within the hour, while the wall's own
temperature is not. That is more easily explained by a response acting on the
instrument or its immediate mounting than by one acting through the depth of the
wall, and it strengthens reading \textbf{(ii)} of Section~\ref{sec:signtest}:
instrument thermal drift dominating a smaller structural response.}

\subsubsection*{How the lead is kept out of every model}

A lead cannot be used by a forecasting system: applying it would require the
probe's future values at the moment the forecast is issued.
\texttt{ip\_lib.operators\_from\_summary} therefore takes a
\texttt{causal\_only} flag, default true, which clamps any negative optimum to
zero and reports that it has done so. Every predictive operator in this report is
built through that guard.

The guard was verified rather than assumed. The thirteen tables carrying every
ranking, subset search, importance diagnostic, out-of-period test and horizon
result were compared byte for byte against the run made before the signed scan
existed. All thirteen are identical. The diagnostic in this subsection is purely
additive: it changes nothing that any model touched.

\subsection{The ranking}

Each driver is then scored on its own, at its diurnal-band operator, by pooled
error over five rolling-origin folds --- every test block strictly after the
training block it was fitted on.

Autoregression is deliberately excluded from this stage. One hour ahead, the
inclination's own recent history explains almost all of its variance, and every
driver added on top of it looks equally negligible; a ranking computed that way
measures how little room the autoregressive term leaves, not what the channel
knows. Autoregression returns in Section~\ref{sec:horizon}, where it belongs.

Battery voltage is carried through as a \textbf{negative control}. It rises with
insolation and temperature, so it correlates with the real drivers, but no
mechanism connects it to the structure.

\begin{center}
\small
\begin{tabular}{lR{1.4cm}R{1.2cm}R{1.6cm}R{1.6cm}R{1.5cm}R{1.9cm}}
\toprule
\textbf{Driver} & \textbf{delay} & \textbf{$\tau$} & \textbf{MAE} &
\textbf{fold sd} & \textbf{$R^2$} & \textbf{available} \\
\midrule
\multicolumn{7}{l}{\itshape Tier 1, four drivers plus the control} \\
\quad \texttt{tair}  & 0 h & 0 h &  8.74 & 7.56 & $+0.760$ & 77.1\,\% \\
\quad \texttt{twall} & 0 h & 0 h & 13.61 & 8.13 & $+0.459$ & 36.6\,\% \\
\quad \texttt{batt}\,$^\dagger$ & 0 h & 1 h & 20.71 & 9.79 & $-0.045$ & 77.1\,\% \\
\quad \texttt{rh}    & 0 h & 0 h & 24.10 & 10.78 & $-0.283$ & 77.1\,\% \\
\quad \texttt{sr}    & 0 h & 1 h & 28.40 & 15.15 & $-0.788$ & 69.1\,\% \\
\addlinespace
\multicolumn{7}{l}{\itshape Tier 2, three drivers plus the control} \\
\quad \texttt{tair}  & 0 h & 0 h &  8.25 & 3.02 & $+0.656$ & 77.1\,\% \\
\quad \texttt{rh}    & 0 h & 0 h & 14.67 & 4.64 & $-0.148$ & 77.1\,\% \\
\quad \texttt{sr}    & 0 h & 1 h & 18.59 & 8.18 & $-1.110$ & 69.1\,\% \\
\quad \texttt{batt}\,$^\dagger$ & 0 h & 1 h & 30.90 & 5.14 & $-3.441$ & 77.1\,\% \\
\bottomrule
\multicolumn{7}{l}{\footnotesize $^\dagger$ negative control. MAE in mdeg;
availability is measured over the whole era, not the window.}
\end{tabular}
\end{center}

\begin{figure}[H]
\includegraphics{S3_F07_st02_rank_tier1}
\caption{Single-predictor skill on the Tier~1 window, without autoregression.
Error bars are the spread across the five rolling-origin folds. The right panel
carries each channel's availability over the whole era, which is what prices wall
temperature's second place.}
\end{figure}

\begin{figure}[H]
\includegraphics{S3_F08_st02_rank_tier2}
\caption{The same ranking on the Tier~2 block. Air temperature keeps first place
with a comparable score and a much tighter fold spread; the ordering below it
changes, which is what a channel whose apparent skill is seasonal rather than
structural looks like.}
\end{figure}

\subsection{What the negative control reveals}

The control ranks \emph{third of five} in Tier~1, ahead of relative humidity and
solar radiation. Read carelessly this looks like a failure of the procedure. Read
correctly it is the procedure working.

The dividing line is not the ordering but the sign of $R^2$. Only \texttt{tair}
($+0.760$) and \texttt{twall} ($+0.459$) beat the mean of the training block.
Battery voltage, relative humidity and solar radiation all have \emph{negative}
out-of-sample $R^2$: each, used alone, is worse than predicting a constant. The
control lands among them, which is exactly what it is for --- it marks the noise
floor, and shows that the three channels below the line carry no standalone
predictive content whatever their in-sample correlation suggests. In Tier~2 the
control falls to last outright.

\panel{Answer to Question 1}{%
\textbf{Air temperature}, applied instantaneously --- delay 0, $\tau = 0$ --- is
the best single predictor in both tiers, with out-of-sample $R^2$ of 0.760 on
Tier~1 and 0.656 on Tier~2 and a mean absolute error of 8.2--8.7\,mdeg. Wall
temperature is the only other channel that beats a constant, at $R^2$ 0.459 once
its implausible 20-hour operator is disallowed. Solar radiation and relative
humidity have no standalone predictive value: both score below a channel with no
physical mechanism.

\medskip
\textbf{No driver requires a transport delay.} Within the twelve-hour physical
window every optimum is instantaneous, and only solar radiation takes a one-hour
time constant.

\medskip
This answer must be read with Section~\ref{sec:signtest} in hand. Air temperature
wins in large part because the target was built by subtracting a multiple of it,
which also explains why the response is instantaneous. The finding that \emph{is}
independent of that construction is the negative one: no channel other than the
two temperatures carries usable information on its own.}

% ===========================================================================
\section{Question 2 --- combinations, and how to diagnose them}
\label{sec:combinations}

\subsection{How far the candidates duplicate one another}
\label{sec:collinearity}

These channels are physically coupled: radiation heats the air, and air and
radiation together heat the wall. A regression containing several of them splits
one physical effect across correlated coefficients in a way that is not stable
between folds.

\begin{center}
\small
\begin{tabular}{lR{1.6cm}R{2.6cm}l}
\toprule
\textbf{Channel} & \textbf{VIF} & \textbf{max $|r|$ with others} &
\textbf{closest} \\
\midrule
\texttt{tair}  & 8.06 & 0.924 & \texttt{twall} \\
\texttt{twall} & 7.90 & 0.924 & \texttt{tair} \\
\texttt{rh}    & 2.69 & 0.732 & \texttt{twall} \\
\texttt{sr}    & 1.82 & 0.610 & \texttt{rh} \\
\bottomrule
\end{tabular}
\end{center}

The two temperatures are near-duplicates at $r = 0.924$, with variance inflation
factors around eight. That is why the search below is scored by out-of-sample
error rather than by coefficients or in-sample fit, and why
Section~\ref{sec:diagnostics} uses two importance measures instead of reading a
coefficient table.

\subsection{Every combination}

With four candidates the search is exhaustive at fifteen subsets, which removes
the need for a stepwise procedure and the selection bias that comes with one.
Each subset is scored three ways: pooled rolling-origin MAE; the Bayesian
information criterion of the corresponding least-squares fit; and an
availability-weighted utility, defined as the ratio of the best subset's error to
this subset's error, multiplied by the fraction of the \emph{whole era} on which
the subset exists.

\begin{center}
\small
\begin{tabular}{lR{0.7cm}R{1.5cm}R{1.4cm}R{1.4cm}R{1.6cm}R{1.4cm}}
\toprule
\textbf{Subset} & \textbf{$k$} & \textbf{MAE} & \textbf{fold sd} &
\textbf{$R^2$} & \textbf{available} & \textbf{utility} \\
\midrule
\multicolumn{7}{l}{\itshape Tier 1 --- top seven of fifteen, by error} \\
\quad \texttt{tair+sr+rh}       & 3 & \textbf{8.24} & 7.15 & 0.782 & 69.1\,\% & 0.691 \\
\quad \texttt{tair+sr}          & 2 & 8.29 & 7.21 & 0.779 & 69.1\,\% & 0.687 \\
\quad \texttt{tair+twall+sr+rh} & 4 & 8.30 & 7.13 & 0.781 & 36.5\,\% & 0.362 \\
\quad \texttt{tair+twall+sr}    & 3 & 8.38 & 7.17 & 0.777 & 36.5\,\% & 0.359 \\
\quad \texttt{tair}             & 1 & 8.74 & 7.56 & 0.760 & 77.1\,\% & \textbf{0.727} \\
\quad \texttt{tair+rh}          & 2 & 8.75 & 7.54 & 0.759 & 77.1\,\% & 0.726 \\
\quad \texttt{tair+twall}       & 2 & 8.80 & 7.53 & 0.759 & 36.6\,\% & 0.343 \\
\addlinespace
\multicolumn{7}{l}{\itshape Tier 2 --- all seven} \\
\quad \texttt{tair+sr}          & 2 & \textbf{5.99} & 0.93 & 0.793 & 69.1\,\% & \textbf{0.691} \\
\quad \texttt{tair+sr+rh}       & 3 & 6.13 & 1.00 & 0.782 & 69.1\,\% & 0.676 \\
\quad \texttt{tair+rh}          & 2 & 8.02 & 3.06 & 0.672 & 77.1\,\% & 0.576 \\
\quad \texttt{tair}             & 1 & 8.25 & 3.02 & 0.656 & 77.1\,\% & 0.560 \\
\quad \texttt{rh}               & 1 & 14.67 & 4.64 & $-0.148$ & 77.1\,\% & 0.315 \\
\quad \texttt{sr+rh}            & 2 & 14.70 & 4.27 & $-0.158$ & 69.1\,\% & 0.282 \\
\quad \texttt{sr}               & 1 & 18.59 & 8.18 & $-1.110$ & 69.1\,\% & 0.223 \\
\bottomrule
\end{tabular}
\end{center}

\panel{Wall temperature is in no winning subset, in either tier}{%
The best Tier~1 subset is \texttt{tair + sr + rh}. The four-channel model that
adds wall temperature is \emph{worse} on error (8.30 against 8.24) and less than
half as available. Every subset containing \texttt{twall} is beaten by the same
subset without it. This is a direct reversal of the earlier, unbounded-delay
version of this study, in which the four-channel model led the table --- and the
reversal is traceable entirely to disallowing a 20-hour delay that the site's
physics excludes.}

\begin{figure}[H]
\includegraphics{S3_F09_st02_subsets_tier1}
\caption{Exhaustive subset search on Tier~1. Left: error against subset size,
coloured by availability, with the best-per-size frontier dashed. The frontier is
almost flat --- adding predictors buys 6\,\% in total. Right: the eight best
subsets once availability is priced in, where the single-predictor model wins
outright.}
\end{figure}

\begin{figure}[H]
\includegraphics{S3_F10_st02_subsets_tier2}
\caption{The same search on Tier~2. Here the frontier has a real step between one
and two predictors, and \texttt{tair + sr} wins on error, on BIC and on utility
simultaneously --- the only configuration in the study on which all three
criteria agree.}
\end{figure}

\subsection{The two tiers give different answers, and only one is resolved}

On Tier~1 the best subset beats the best single predictor by 0.50\,mdeg, that is
6.0\,\%. The fold-to-fold spread of both estimates is around 7.2--7.6\,mdeg,
fifteen times the difference between them. \textbf{The Tier~1 ordering below the
top place is not resolved by these data.} What Tier~1 does establish is that no
combination is dramatically better than air temperature alone, and that on
utility the single predictor wins outright.

On Tier~2 the picture is different and much cleaner. Adding solar radiation to
air temperature reduces the error from 8.25 to 5.99\,mdeg, that is
\textbf{27.4\,\%}, and lifts $R^2$ from 0.656 to 0.793. The fold spread of the
winning model is 0.93\,mdeg, well inside the 2.26\,mdeg improvement, although the
single-predictor model's own spread of 3.02\,mdeg means the comparison is
directionally clear rather than sharply resolved. BIC agrees, dropping by 921,
and utility agrees despite \texttt{sr} costing 8 percentage points of
availability.

Why the tiers disagree is worth stating: Tier~2 runs from May to September and
contains the high-summer period in which radiation is both strong and variable,
while Tier~1 is dominated by spring, when the radiation signal is weaker relative
to the air-temperature signal that already dominates the target. Solar radiation
earns its place in the season in which it is doing work.

\subsection{Which member is doing the work}
\label{sec:diagnostics}

Two diagnostics are run on the winning Tier~1 subset, deliberately different,
because their disagreement is informative.

\textbf{Leave-one-out gain} removes a channel and \emph{refits}. It answers: what
would it cost to build this system without that channel? Under strong
collinearity a refit can recover much of the lost information from a correlated
channel, so this number is small exactly when the channel is replaceable.

\textbf{Permutation importance} shuffles a channel within the test block while
holding the fitted model fixed. It answers: how much does the model \emph{as
built} rely on that channel?

\begin{center}
\small
\begin{tabular}{lR{2.0cm}R{1.8cm}R{2.6cm}R{1.4cm}}
\toprule
\textbf{Channel} & \textbf{MAE without} & \textbf{gain} & \textbf{permutation
increase} & \textbf{sd} \\
\midrule
\texttt{tair} & 24.80 & 16.56 (66.8\,\%) & 8.84 & 2.40 \\
\texttt{sr}   &  8.75 &  0.51 (5.9\,\%)  & 0.013 & 0.019 \\
\texttt{rh}   &  8.29 &  0.05 (0.6\,\%)  & $-0.010$ & 0.026 \\
\bottomrule
\multicolumn{5}{l}{\footnotesize Subset \texttt{tair + sr + rh}; full-subset MAE
8.24\,mdeg. All values in mdeg.}
\end{tabular}
\end{center}

\begin{figure}[H]
\includegraphics{S3_F11_st02_diagnostics}
\caption{Contribution of each channel to the winning Tier~1 subset. Left: the
cost of removing the channel and refitting. Right: the cost of shuffling it with
the model held fixed, with the spread over folds and repeats. Air temperature
dominates both. Both other channels sit at or below their own error bar on the
right-hand panel, which is the quantitative statement that the model does not
rely on them.}
\end{figure}

Air temperature dominates both measures by three orders of magnitude on the
permutation scale. Solar radiation and relative humidity each produce a
permutation increase smaller than its own standard deviation, so the fitted model
does not measurably rely on either. \texttt{sr} shows the one interesting
asymmetry --- a 5.9\,\% leave-one-out gain against a permutation increase of
essentially zero --- but at 0.51\,mdeg against a 7.15\,mdeg fold spread, that
asymmetry is not resolved either.

\subsection{The out-of-period test}

The Tier~1 models are refitted on the whole Tier~1 window and applied to the
detached 2026 summer block, fourteen months later, with no refitting. This is the
hardest test in the study: a different season, a different part of the annual
cycle, and an instrument that has been through a 103-day outage in between.

\begin{center}
\small
\begin{tabular}{lR{1.4cm}R{1.6cm}R{1.6cm}}
\toprule
\textbf{Model} & \textbf{$n$} & \textbf{MAE} & \textbf{RMSE} \\
\midrule
\texttt{tair + sr + rh} & 1225 & 17.39 & 20.08 \\
\texttt{tair} alone     & 1225 & \textbf{16.96} & 20.12 \\
\bottomrule
\end{tabular}
\end{center}

Both models roughly double their cross-validated error, and \textbf{the single
predictor is marginally the better of the two}. Whatever the two additional
channels contributed inside the Tier~1 window did not survive the transfer. This
is the strongest single piece of evidence in the section, and it points the same
way as the utility column.

\subsection{The diagnostic protocol}

Question~2 asks not only which combination wins but how to tell. The procedure
used above is the answer, and it transfers to any station and any sensor set.

\begin{center}
\small
\begin{tabularx}{\textwidth}{R{0.5cm}L{6.0cm}Y}
\toprule
& \textbf{Rule} & \textbf{Failure it prevents} \\
\midrule
1 & Bound the operator scan by the physics of the mechanism, and by half the
dominant period &
A physically impossible optimum, or an aliased lead reported as a delay \\
2 & Fit each driver's operator first, on the band of interest &
A lagged driver ranked below an instantaneous one; a boundary optimum reported as
a physical time constant \\
3 & Score without autoregression &
Every driver tying behind the autoregressive term at short horizons \\
4 & Use rolling-origin folds and report the spread &
A ranking that depends on where a single split happened to fall \\
5 & Search subsets exhaustively when the candidate set is small &
Selection bias and instability under collinearity \\
6 & Never read coefficients under high VIF; use leave-one-out gain and
permutation importance, and read their disagreement as redundancy &
Importance attributed by collinearity rather than by content \\
7 & Price availability over the whole record, not over the window &
Choosing a channel that is absent when it is needed \\
8 & Carry a negative control &
Mistaking shared diurnal shape for predictive content \\
9 & Confirm on an out-of-period block &
A combination that fits one season and transfers to none \\
\bottomrule
\end{tabularx}
\end{center}

Rule 1 is first because it is the rule this study learned the hard way: the
earlier unbounded scan changed the identity of the winning subset, and no
downstream diagnostic would have caught it.

\panel{Answer to Question 2}{%
\textbf{Yes, in the season where the extra channel does work --- and never with
wall temperature.} On the Tier~2 summer block \texttt{tair + sr} beats
\texttt{tair} alone by 27.4\,\% in mean absolute error, and wins on error, BIC
and availability-weighted utility simultaneously. On the spring-dominated Tier~1
window the best subset improves on air temperature by 6.0\,\%, a fifteenth of the
fold-to-fold spread, and the out-of-period test at fourteen months prefers the
single predictor outright.

\medskip
\textbf{Wall temperature appears in no winning subset in either tier}, as a
\emph{predictor}. Every subset containing it is beaten by the same subset without
it, at less than half the availability. On this evidence it should not be a
required regressor of the production pipeline --- which resolves, in the
negative, the open question raised in
\texttt{docs/data-quality-report-2026-08-10.md}, Section~6.

\medskip
\textbf{That verdict is about forecasting, and must not be read as ``the channel
is uninformative''.} Section~\ref{sec:probe} shows the opposite: given a signed
operator, \texttt{twall} explains $R^2 = 0.797$ of the diurnal band against 0.491
under the causal constraint, and it is the only channel in the package that loses
anything at all to that constraint. It fails as a predictor precisely
\emph{because} what it knows arrives too late to forecast with --- which is
itself a measurement of where the probe sits. A channel can be worthless for
prediction and valuable for diagnosis, and this one is both.

\medskip
The diagnosis is made by the nine-rule protocol above; the four components that
did the actual work here were the physical bound on the operator scan, the
availability-weighted utility, the disagreement between the two importance
measures, and the out-of-period block.}

% ===========================================================================
\section{Question 3 --- how far ahead, and with what uncertainty}
\label{sec:horizon}

\subsection{Design}

NeuralProphet is fitted with \texttt{n\_lags} $=24$ and \texttt{n\_forecasts}
$=168$, so a single training run yields every horizon from one hour to one week
and the whole sweep comes from one fit rather than from eight. Quantile
regression at the 5th and 95th percentiles attaches a nominal 90\,\% prediction
interval to every forecast. Gaps are bridged by interpolation so that the
autoregressive window has continuous history, and scoring happens only at
originally observed timestamps: the model may read a filled value, but it is
never rewarded for reproducing one.

Two regressor regimes are run, and the difference between them is the point of
the exercise:

\begin{itemize}
\item \textbf{lagged} --- the regressors enter only through their past values.
This is what a deployed system has, and it is the honest number.
\item \textbf{future} --- the regressors are known across the forecast window.
This is what the same system coupled to a \emph{perfect} weather forecast would
have. On this target it is more than that, and the excess is discussed in
Section~\ref{sec:regimegap}.
\end{itemize}

Skill is measured against two baselines. \textbf{Persistence} carries the last
observation forward. \textbf{Seasonal naive} repeats the value from a whole
number of days earlier, which is the cheapest possible way to exploit the diurnal
cycle. At horizons that are exact multiples of 24 hours the two baselines
coincide by construction.

\subsection{Results}

\begin{center}
\small
\begin{tabular}{R{1.3cm}R{1.5cm}R{1.5cm}R{1.9cm}R{2.0cm}R{1.6cm}R{1.6cm}}
\toprule
\textbf{$h$ [h]} & \textbf{MAE} & \textbf{persist.} & \textbf{seas.\ naive} &
\textbf{skill vs p.} & \textbf{cover.} & \textbf{width} \\
\midrule
\multicolumn{7}{l}{\itshape Tier 2, lagged --- the deployable configuration} \\
1   &  4.43 &  2.79 & 6.10 & $-0.585$ & 61.4\,\% & 15.3 \\
3   &  5.59 &  6.81 & 6.09 & $+0.179$ & 54.4\,\% & 16.3 \\
6   &  6.35 & 11.74 & 6.09 & $+0.459$ & 56.6\,\% & 24.6 \\
12  &  6.35 & 16.06 & 6.10 & $\mathbf{+0.605}$ & 53.6\,\% & 18.4 \\
24  &  6.96 &  6.10 & 6.10 & $-0.141$ & 57.1\,\% & 22.9 \\
48  &  8.09 &  8.30 & 8.30 & $+0.025$ & 67.1\,\% & 27.7 \\
72  &  9.62 &  9.67 & 9.67 & $+0.005$ & 66.8\,\% & 31.7 \\
168 & 15.95 & 12.22 & 12.22 & $-0.306$ & 58.1\,\% & 29.7 \\
\addlinespace
\multicolumn{7}{l}{\itshape Tier 2, future --- regressors known over the horizon} \\
1   &  3.82 &  2.86 & 6.13 & $-0.336$ & 95.4\,\% & 40.0 \\
6   &  5.88 & 12.05 & 6.11 & $+0.512$ & 95.2\,\% & 42.8 \\
12  &  6.07 & 16.62 & 6.13 & $+0.635$ & 95.8\,\% & 45.5 \\
24  &  6.50 &  6.12 & 6.12 & $-0.061$ & 96.7\,\% & 47.9 \\
72  &  9.06 &  9.49 & 9.49 & $+0.046$ & 93.0\,\% & 53.1 \\
168 & 13.07 & 12.84 & 12.84 & $-0.018$ & 82.7\,\% & 50.6 \\
\addlinespace
\multicolumn{7}{l}{\itshape Tier 1, future --- accurate, but see
Section~\ref{sec:regimegap}} \\
1   &  2.03 &  3.42 & 5.33 & $+0.406$ & 18.0\,\% & 0.0 \\
12  &  2.47 & 21.61 & 5.33 & $+0.886$ & 18.9\,\% & 0.0 \\
24  &  2.85 &  5.33 & 5.33 & $+0.464$ & 23.4\,\% & 0.0 \\
168 &  7.36 & 10.96 & 10.96 & $+0.328$ & 21.3\,\% & 0.0 \\
\addlinespace
\multicolumn{7}{l}{\itshape Tier 1, lagged --- fitted across the 5 May
discontinuity} \\
1   &  4.43 &  3.44 & 5.76 & $-0.285$ & 12.3\,\% & 0.0 \\
12  & 14.42 & 21.20 & 5.76 & $+0.320$ &  9.8\,\% & 0.0 \\
24  & 21.67 &  5.76 & 5.76 & $-2.764$ & 11.2\,\% & 0.2 \\
168 & 52.82 & 10.62 & 10.62 & $-3.974$ &  0.0\,\% & 2.7 \\
\bottomrule
\multicolumn{7}{l}{\footnotesize MAE and interval width in mdeg. Coverage is
against a nominal 90\,\%.}
\end{tabular}
\end{center}

\begin{figure}[H]
\includegraphics{S3_F12_st02_horizon_skill}
\caption{Forecast error (left) and skill against persistence (right) for all four
configurations. The dashed line in the left panel is persistence itself, and its
shape explains most of the right panel: persistence is excellent at one hour,
worst at twelve, and excellent again at twenty-four because the signal is
diurnal. Skill against it is therefore not monotonic in the horizon, and quoting
a single ``maximum useful horizon'' would misdescribe the curve.}
\end{figure}

\subsection{Reading the horizon curve}

The non-monotonic skill curve is the substantive finding, and it follows from the
structure of the signal rather than from any property of the model. Taking the
deployable Tier~2 lagged configuration:

At \textbf{one hour} nothing beats persistence. The inclination changes slowly
enough that its own last value is an excellent forecast (MAE 2.79\,mdeg), and no
configuration in the study --- including a model trained for that horizon alone
--- gets below it.

Between \textbf{three and twelve hours} the model has real skill, peaking at
$+0.605$ at twelve hours. This is the interval in which persistence has decayed
but the diurnal cycle has not yet brought the signal back, and it is where the
environmental regressors earn their place.

At \textbf{twenty-four hours} the diurnal periodicity restores both baselines to
6.10\,mdeg and the model, at 6.96, is beaten by simply repeating yesterday.

Beyond \textbf{forty-eight hours} the model and the baselines converge, with
skill of $+0.025$ and $+0.005$ --- indistinguishable from zero. At one week the
model is 30\,\% worse than repeating the value from a week before.

The Tier~1 lagged configuration fails outright: 21.7\,mdeg at twenty-four hours
and 52.8 at one week, against baselines of 5.8 and 10.6. The diagnosis is in
Section~2.2 --- the Tier~1 window is two runs separated by a wall-probe
interruption on 5 May, so an autoregressive model fitted on it is fitted across a
discontinuity that its lag window cannot see. This is not a finding about wall
temperature; it is a finding about training an autoregressive model on a
fragmented record, and it is why Tier~2 is the configuration of record.

\subsection{The gap between the regimes is not the value of a weather forecast}
\label{sec:regimegap}

The Tier~1 future configuration is the most accurate in the study by a wide
margin: 2.03\,mdeg at one hour, 2.47 at twelve, 7.36 at one week, with positive
skill against persistence at \emph{every} horizon. Taken at face value this says
that a system coupled to a perfect weather forecast could predict this structure
a week ahead.

It should not be taken at face value. The target contains $-5\,T_{\mathrm{air}}$
by construction (Section~\ref{sec:signtest}), so handing the model the future air
temperature is close to handing it the future target. The future regime is
therefore measuring two things at once --- the genuine value of knowing the
weather, and the arithmetic value of being told the term that was subtracted ---
and on this target the second dominates.

The corollary is that the lagged--future gap in Tier~2, from 4.43 to 3.82\,mdeg
at one hour and 6.35 to 6.07 at twelve, is an \emph{upper} bound on what a
weather forecast is worth here, and even that bound is modest: 14\,\% and 4\,\%
respectively. The main benefit of the future regime is not accuracy but
calibration, discussed next.

\subsection{Is the quoted uncertainty honest?}

\begin{figure}[H]
\includegraphics{S3_F13_st02_uncertainty}
\caption{Calibration (left) and width (right) of the nominal 90\,\% prediction
interval. The dashed line marks the nominal rate. Only the Tier~2 future
configuration tracks it. The deployable Tier~2 lagged configuration sits 23 to 36
percentage points below while producing intervals roughly half as wide --- the
signature of an overconfident model. Both Tier~1 configurations return degenerate
intervals of zero width and are plotted only to document the failure.}
\end{figure}

This is the part of Question~3 with the most operational consequence.

The \textbf{Tier~2 future} configuration is well calibrated: empirical coverage
of 93--97\,\% against a nominal 90\,\% out to three days, falling to 82.7\,\% at
one week. Slightly conservative, and usable as quoted.

The \textbf{Tier~2 lagged} configuration --- the one that can actually be
deployed --- covers only 53.6--67.1\,\%. Its nominal ninety per cent interval is
in reality a fifty-five to sixty-five per cent interval. Its median width, 15 to
32\,mdeg, is between 38 and 60\,\% of the future-regime width, so the model is
not merely mis-stating a correct interval but is confidently narrow about a
prediction it cannot support.

\textbf{Both Tier~1 configurations produce degenerate quantile heads}, with
interval widths of exactly zero at almost every horizon and coverage between 0
and 23\,\%. The Tier~1 future configuration is therefore accurate and
\emph{completely uninformative about its own uncertainty} --- a combination worth
naming, because a point forecast that good invites trust that its interval does
not earn.

\panel{Uncertainty must be recalibrated before it is published}{%
The interval NeuralProphet produces under the deployable regime is not a 90\,\%
interval. Any downstream anomaly threshold built on it would fire on three to
four times the intended fraction of ordinary observations: the intended miss rate
is 10\,\%, the measured one runs from 33 to 46\,\%. Two routes are open and both
are outside this study's scope: split-conformal recalibration on a held-out
block, which would restore nominal coverage by construction; or widening the
quantile spread until measured coverage reaches ninety, which is the same
operation done by hand. Until one of them is applied, only the measured coverage
should be quoted, never the nominal one.}

\begin{figure}[H]
\includegraphics{S3_F14_st02_fan}
\caption{Ten days of the Tier~2 lagged test period at a twenty-four-hour horizon,
with the 90\,\% interval. The forecast reproduces the diurnal shape and the broad
level, and the interval is visibly too narrow for the excursions it is meant to
contain --- the picture behind the 57\,\% coverage in the table above.}
\end{figure}

\subsection{What the joint multi-horizon head costs}

One fit produced every horizon, which is efficient but not free: the network
shares one set of weights across 168 output steps. Refitting with
\texttt{n\_forecasts} set to exactly one horizon isolates the cost.

\begin{center}
\small
\begin{tabular}{R{1.6cm}R{2.2cm}R{2.2cm}R{2.4cm}}
\toprule
\textbf{$h$ [h]} & \textbf{dedicated} & \textbf{joint head} &
\textbf{joint penalty} \\
\midrule
1  & 3.30 & 4.43 & $+34.2$\,\% \\
24 & 7.45 & 6.96 & $-6.6$\,\% \\
\bottomrule
\end{tabular}
\end{center}

The joint head costs 34\,\% at one hour and \emph{gains} 6.6\,\% at twenty-four,
consistent with the shared weights acting as a regulariser at longer horizons and
as a distraction at the shortest one. The conclusion is unchanged either way:
even the dedicated one-hour model, at 3.30\,mdeg, does not beat persistence at
2.79.

\panel{Answer to Question 3}{%
\textbf{The useful forecast band is roughly three to twelve hours ahead}, where
skill against persistence runs from $+0.18$ to $+0.61$ and the mean absolute
error is 5.6--6.4\,mdeg. Below three hours persistence wins; at twenty-four hours
and beyond, repeating the previous day wins. There is no horizon at which the
deployable system beats a naive baseline by more than a factor of two and a half.

\medskip
\textbf{The uncertainty attached to that band is not what it claims.} The
deployable configuration's nominal 90\,\% interval covers 54--67\,\% of
observations, and must be recalibrated before publication. Only the Tier~2 future
configuration is calibrated as quoted.

\medskip
\textbf{A weather forecast is worth having but not transformative}: at most
14\,\% at one hour and 4\,\% at twelve, and its principal benefit is calibration
rather than accuracy. The far larger apparent gain in Tier~1 is an artefact of
the target's construction, not a forecast benefit.}

% ===========================================================================
\section{Verdicts}

\begin{center}
\small
\begin{tabularx}{\textwidth}{L{4.4cm}Y}
\toprule
\textbf{Question} & \textbf{Answer} \\
\midrule
0 · expected sign &
Site physics predicts negative; the compensated target complies at $r = -0.956$,
the \textbf{raw channel does not} ($r = +0.845$, $+1.63$\,\mdegC). The compliance
is arithmetic, not structural. \fail\ as evidence for the mechanism. \\
1 · best single predictor &
\texttt{tair}, instantaneous, in both tiers. $R^2$ 0.760 (Tier~1) and 0.656
(Tier~2); MAE 8.74 and 8.25\,mdeg. \\
1 · lag structure &
No \emph{forcing} needs a transport delay within the twelve-hour physical window;
only \texttt{sr} takes a time constant, of one hour. \\
1 · wall-probe phase &
The inclination \textbf{leads} \texttt{twall} by 2\,h, on the raw channel as well
as the compensated one, and the signed optimum is $-4$\,h at $R^2$ 0.797 against
0.491 causal. Signature of a \textbf{deep} probe. Depth unrecorded; the phase is
the check to run once it is known. \\
1 · what else is usable &
Only \texttt{twall} ($R^2$ 0.459) also beats a constant. \texttt{sr}, \texttt{rh}
and the negative control all score below a constant. \\
2 · best combination &
\texttt{tair + sr} on Tier~2: $-27.4$\,\% MAE, best on error, BIC and utility
together. On Tier~1 the 6.0\,\% gain is not resolved above the fold spread. \\
2 · wall temperature &
\fail\ \emph{as a predictor}. Present in no winning subset in either tier; every
subset containing it is beaten by the same subset without it, at 36.6\,\%
availability. \pass\ \emph{as a diagnostic} --- see the wall-probe phase row. \\
2 · negative control &
Third of five in Tier~1, last in Tier~2 --- and the boundary between positive and
negative out-of-sample $R^2$ falls above it. \\
3 · useful horizon &
3 to 12 hours. Peak skill $+0.605$ at 12\,h. Persistence wins at 1\,h; the
previous day wins at 24\,h and beyond. \\
3 · uncertainty &
\fail\ as quoted. Nominal 90\,\% interval measures 54--67\,\% under the
deployable regime; both Tier~1 configurations are degenerate. \\
\bottomrule
\end{tabularx}
\end{center}

% ===========================================================================
\section{Caveats, and what would change the answers}

\textbf{The target is partly artificial, and the sign test proves it.} The raw
channel moves the other way from the compensated one. The single change most
likely to alter every conclusion here is re-running the study on the raw
inclination, or on a compensated series using the 1.63\,\mdegC\ slope the
companion study fitted. The methods transfer unchanged; only the target moves.

\textbf{The expected sign is unresolved, not confirmed.} Section~\ref{sec:signtest}
lists three readings consistent with the measurement, and the archive contains
nothing that distinguishes them. The calibration sheet or the installation record
would.

\textbf{The delay bound is an assumption, stated and defended.} Twelve hours
follows from the mechanism and from half the diurnal period. It changed the
winning subset, so it is not a neutral choice --- but for an external forcing the
alternative admits optima that the site's physics excludes.

\textbf{The wall probe's depth is unknown, and it is the cheapest open question
to close.} It decides whether the two-hour lead of
Section~\ref{sec:probe} is a shallow or a deep installation, and it bears
directly on the sign contradiction of Section~\ref{sec:signtest}. The phase is on
record, so answering it requires no re-analysis --- only a comparison. It is
listed as open question 3 in \texttt{docs/data-quality-report-2026-08-10.md}.

\textbf{The tiers overlap and are short.} Tier~1 and Tier~2 share 70 days.
Neither exceeds 144 days, so no annual term is identifiable and every seasonal
statement here is an inference from two seasons rather than a measurement over a
cycle.

\textbf{Tier~1 is internally fragmented.} The 5 May interruption splits it, which
is enough to break the autoregressive model in Section~\ref{sec:horizon}. A
production system must screen for this before fitting, not after.

\textbf{Ridge is a screening model, not the deployed one.} It orders candidate
sets cheaply so that NeuralProphet is run only on the survivors. Where the two
disagree, NeuralProphet is the relevant answer.

\textbf{The horizon study uses one chronological split per configuration.} The
70\,\% cut is the same for all four, so the comparison between them is sound, but
the absolute skill numbers carry the uncertainty of a single origin.

% ===========================================================================
\section{Artefact inventory}

Every figure and table in this report is produced by
\texttt{inclination\_prediction\_study.ipynb} and written to \texttt{outputs/}.

\begin{center}
\small
\begin{tabularx}{\textwidth}{L{4.6cm}Y}
\toprule
\textbf{Artefact} & \textbf{Content} \\
\midrule
\texttt{S3\_01} -- \texttt{S3\_04} & Channel summary, coverage, contiguous
blocks with and without wall temperature \\
\texttt{S3\_05}, \texttt{S3\_06} & Correlation matrices, levels and differences \\
\texttt{S3\_23} & Sign test: raw against compensated, correlation and slope \\
\texttt{S3\_07}, \texttt{S3\_08}, \texttt{S3\_10} & Operator optima: full band,
diurnal band, Tier~2 \\
\texttt{S3\_24}, \texttt{S3\_25} & Signed operator optima and the probe-phase
table --- diagnostic, feeding no model \\
\texttt{S3\_09}, \texttt{S3\_11} & Single-predictor rankings, both tiers \\
\texttt{S3\_12} & Variance inflation factors \\
\texttt{S3\_13}, \texttt{S3\_14} & Exhaustive subset searches, both tiers \\
\texttt{S3\_15}, \texttt{S3\_16} & Leave-one-out gain, permutation importance \\
\texttt{S3\_17} & Out-of-period test on the 2026 block \\
\texttt{S3\_18} & The diagnostic protocol \\
\texttt{S3\_19} & Horizon tables, four configurations \\
\texttt{S3\_20} & Horizon limits and coverage summary \\
\texttt{S3\_21} & Dedicated against joint multi-horizon head \\
\texttt{S3\_22} & Verdicts \\
\addlinespace
\texttt{S3\_F01} -- \texttt{S3\_F04} & Overview, target, diurnal cycle,
correlations \\
\texttt{S3\_F05}, \texttt{S3\_F06} & Operator heat maps, both bands \\
\texttt{S3\_F07}, \texttt{S3\_F08} & Single-predictor rankings \\
\texttt{S3\_F09}, \texttt{S3\_F10} & Subset searches \\
\texttt{S3\_F11} & Contribution diagnostics \\
\texttt{S3\_F15} & Signed-delay operator grid, probe-phase diagnostic \\
\texttt{S3\_F12}, \texttt{S3\_F13}, \texttt{S3\_F14} & Horizon skill,
calibration, fan chart \\
\bottomrule
\end{tabularx}
\end{center}

Reproduction, from the study directory with the \texttt{neuralprophet\_env}
environment active:

\begin{quote}
\ttfamily\small
jupytext -{}-to ipynb inclination\_prediction\_study.py \\
jupyter nbconvert -{}-to notebook -{}-execute -{}-inplace \\
\phantom{jupyter nbconvert }inclination\_prediction\_study.ipynb
\end{quote}

Figures are drawn through seaborn; setting \texttt{CONTEXT = 'paper'} in the
parameter cell re-exports every one of them at manuscript column size. The
site's sign convention and the contradiction of Section~\ref{sec:signtest} are
recorded durably in \texttt{docs/raw-data-format.md} Section~7.5.

\end{document}
